\documentclass[11pt,a4paper]{article}
\usepackage[margin=2.5cm]{geometry}
\usepackage{amsmath,amssymb,amsthm}
\usepackage{enumitem}
\usepackage[hidelinks]{hyperref}
\usepackage[switch]{lineno}
\setlength{\emergencystretch}{2em}
\newtheorem{theorem}{Theorem}[section]
\newtheorem{lemma}[theorem]{Lemma}
\newtheorem{proposition}[theorem]{Proposition}
\newtheorem{corollary}[theorem]{Corollary}
\newtheorem{definition}[theorem]{Definition}
\newcommand{\D}{\mathcal{D}}
\newcommand{\WitnessSet}{\mathsf{W}}
\newcommand{\cf}{\mathrm{cf}}
\newcommand{\lin}{\mathrm{lin}}
\newcommand{\Ccf}[1]{\mathcal C_{#1}^{\cf}}
\newcommand{\Clin}[1]{\mathcal C_{#1}^{\lin}}
\newcommand{\pref}{\mathrm{pref}}
\newcommand{\suf}{\mathrm{suf}}
\newcommand{\keywords}[1]{\smallskip\noindent{\bfseries keywords:} #1}
\title{Distributional Learning of Context-Free Languages under Fixed Finite-Monoid Typing}
\author{Takayuki Kuriyama\\
Independent Researcher\\
\texttt{growup.kuriyama@gmail.com}}
\date{}
\begin{document}
\linenumbers
\maketitle

\begin{abstract}
We study distributional learning of context-free languages under a fixed recognizable congruence $\sim_h$, the kernel of a homomorphism $h:\Sigma^*\to M$ into a finite monoid.  A yield-typed refinement records the $h$-type of each nonterminal yield.  We define a set-driven grammar reconstruction operator on arbitrary finite positive samples and show that finitely many canonical positive witnesses suffice for exact reconstruction of every context-free $\sim_h$-substitutable target.  Gold identification is obtained separately by a conservative sequential learner that changes its hypothesis only when a new positive example is not generated by the current grammar; for fixed $h$, every reconstruction and learner update is polynomial in the accumulated sample size.

For every fixed prefix--suffix window $(k,\ell)$, including $(0,0)$, the induced finite-monoid typing exactly captures Yoshinaka's $(k,\ell)$-substitutability class.  On this special case, the present reconstruction recovers the same grammar-size-and-thickness characteristic-data scale established by Yoshinaka via a polynomial thickness-preserving normalization to the start-separated binary normal form used in the reconstruction theorem.  For linear targets under an arbitrary fixed $h$, the one-spine structure directly bounds canonical yields and canonical contexts polynomially in the size of a linear CFG representation, yielding polynomial characteristic data without a separate thickness parameter in the statement.  The fixed-window polynomial depends on the fixed window and alphabet; no uniform polynomial bound is claimed when $k,\ell$ vary.  We also give nonregular linear and nonlinear targets showing that fixed finite-monoid comparison reaches beyond every fixed window while still excluding some deterministic context-free languages already over a two-letter alphabet.

\keywords{grammatical inference, distributional learning, context-free languages, recognizable congruences, finite monoid homomorphisms, identification in the limit, exact reconstruction}
\end{abstract}

\section{Introduction}

\subsection{Background and motivation}

Gold~\cite{gold1967} showed that no language class containing all finite
languages and at least one infinite language is identifiable in the limit
from positive data alone. This negative result motivated a broad line of
research in grammatical inference; see, for example,
de~la~Higuera~\cite{higuera2010} for general background. For context-free languages, one standard way to
recover positive learning results is the \emph{distributional} approach,
in which learnability is obtained by restricting which substrings may be
compared across contexts. Classical examples include Clark--Eyraud
substitutability~\cite{clark-eyraud2007} and Yoshinaka's
$(k,\ell)$-substitutability~\cite{yoshinaka2008}. More generally,
positive-data learnability has also been established for other grammar
classes under suitable finite control assumptions; see, for example,
Kanazawa~\cite{kanazawa1998} for restricted categorial grammars.

The present paper studies a recognizable-congruence version of this
distributional program. Let $h:\Sigma^*\to M$ be a homomorphism into a
finite monoid, and write $\sim_h:=\ker h$, that is, $x\sim_h y$ iff $h(x)=h(y)$. For $x\in\Sigma^*$, let
$\D_L(x):=\{(u,v)\in\Sigma^*\times\Sigma^*:uxv\in L\}$ denote its
distribution in $L$. A language $L\subseteq\Sigma^*$ is
\emph{$\sim_h$-substitutable} if $h(x)=h(y)$ and
$\D_L(x)\cap\D_L(y)\neq\emptyset$ imply $\D_L(x)=\D_L(y)$ for all
$x,y\in\Sigma^+$%
\footnote{Clark--Eyraud treat $a^+$ as substitutable, whereas including $\lambda$ itself among the compared fragments would exclude $a^+$. See also Clark~\cite{clark2013}.}. Thus full equality of distributions is required only
for strings that already agree under a fixed finite algebraic summary.

Recognizable congruences are a natural comparison mechanism here. If
$L=h^{-1}(\mathsf{Acc})$ for some finite monoid homomorphism $h:\Sigma^*\to M$ and
some $\mathsf{Acc}\subseteq M$, then any two strings with the same $h$-value
automatically have the same continuation behavior, and hence the same
distribution. In particular, every regular language belongs to the broader
class $\mathsf{RS}$ (defined formally in Section~\ref{sec:framework}) of
languages that are $\sim_h$-substitutable for \emph{some} finite monoid
homomorphism. This observation shows that finite
monoid typings are not an ad hoc generalization of bounded contexts, but a
canonical algebraic source of a finite comparison bias.

Our focus is the \emph{fixed-$h$} regime: the same finite monoid
homomorphism $h:\Sigma^*\to M$ is fixed uniformly over the class
$\Ccf{h}$ of context-free $\sim_h$-substitutable languages, rather than
selected separately for each target after observing its data.  Thus fixed $h$
plays the same class-level role as a fixed pair $(k,\ell)$ in Yoshinaka's
hierarchy.  We separate choosing a useful typing
from learning once that typing has been specified, a viewpoint close to
typing or domain bias in automata inference~\cite{coste-typing2004}.

There is also a useful contrast with control-set approaches.  Takada's
even-linear framework controls derivations of a fixed universal grammar by a
regular language of rule labels~\cite{takada1988,takada1995}, whereas the
present homomorphism acts on terminal yields and supplies a finite
\emph{comparison bias}: fixing $h$ restricts which substrings may be identified
but does not determine one target language.  Related positive-data work on
control sets imposes additional restrictions~\cite{koshiba-makinen-takada1997}.
We use this contrast only for positioning; no class separation from the full
control-set framework is claimed.

\subsection{Contributions and main results}

The paper has three main contributions.

\begin{enumerate}[label=\textup{(\arabic*)},leftmargin=2.4em]
\item For a fixed finite monoid homomorphism $h$, we define a set-driven
reconstruction operator and prove exact reconstruction from finitely many
canonical positive witnesses.  A separate conservative sequential learner then
gives Gold identification for every nonempty language in $\Ccf{h}$, with
polynomial-time reconstruction and updates.

\item For every fixed prefix--suffix window $(k,\ell)$, the induced typing is
exactly equivalent to Yoshinaka's $(k,\ell)$-substitutability condition, and the
general reconstruction recovers the same grammar-size-and-thickness
characteristic-data scale established by Yoshinaka after a polynomial thickness-preserving normalization to
SSBNF.  For the linear subclass under an arbitrary fixed $h$, we instead obtain
direct polynomial bounds on canonical yields and canonical contexts, yielding
characteristic data polynomial in the size of a linear CFG representation.

\item We locate the framework structurally.  The fixed-window hierarchy is
strictly contained in recognizable-congruence substitutability already at the
regular level, and the separation persists for nonregular linear and nonlinear
context-free languages.  Conversely, natural deterministic context-free
languages such as the one-bracket Dyck language lie outside $\mathsf{RS}$.
The class $\mathsf{RS}\cap\mathrm{CFL}$ is contained in Clark's
congruential family, and this containment is proper already over the
two-letter alphabet $\{a,b\}$.
\end{enumerate}

\subsection{Organization of the paper}
Section~\ref{sec:prelim} fixes notation and the learning model, and
Section~\ref{sec:framework} relates fixed recognizable substitutability to
the classical classes.  Section~\ref{sec:algo} defines the reconstruction
operator and the conservative sequential learner, while
Section~\ref{sec:typing} proves exact reconstruction and Gold convergence by
yield-typing a target grammar.  Section~\ref{sec:complexity} gives the
polynomial construction bound.  Section~\ref{sec:fixed-window-thick} verifies that the reconstruction
recovers Yoshinaka's fixed-window grammar-size-and-thickness scale using a
polynomial thickness-preserving SSBNF normalization, and Section~\ref{sec:linear}
proves direct polynomial bounds on typed witnesses for arbitrary fixed-$h$ linear
targets.  Section~\ref{sec:boundaries}
records the main expressiveness and boundary results.  The technical
normalization and short-witness proofs are deferred to the appendices.

\section{Preliminaries}
\label{sec:prelim}

Let $\Sigma$ be a finite alphabet, let $\Sigma^*$ be the free monoid with
identity $\lambda$, and put $\Sigma^+:=\Sigma^*\setminus\{\lambda\}$.
A pair $(u,v)\in\Sigma^*\times\Sigma^*$ is called a
\emph{terminal context}.  For $L\subseteq\Sigma^*$ and $x\in\Sigma^*$, the \emph{distribution} of $x$ in $L$ is
$\D_L(x):=\{(u,v)\in\Sigma^*\times\Sigma^*:uxv\in L\}$.
The \emph{syntactic congruence} of $L$ is defined by $x\equiv_L y$ iff
$\D_L(x)=\D_L(y)$, for $x,y\in\Sigma^*$.
We say that $x$ and $y$ are \emph{syntactically congruent} when
$x\equiv_L y$, and write $[x]_{\equiv_L}$ for the $\equiv_L$-equivalence
class of $x$.

A context-free grammar (CFG) is a tuple $G=(V,\Sigma,P,S)$ with finite
nonterminal set $V$, start symbol $S\in V$, and productions
$A\to\alpha$ with $A\in V$ and $\alpha\in(V\cup\Sigma)^*$.
A \emph{sentential form} is any word in $(V\cup\Sigma)^*$.  We write
$\alpha\Rightarrow_G\beta$ for one derivation step in $G$, and
$\Rightarrow_G^*$ for its reflexive-transitive closure.  Thus $L(G):=\{w\in\Sigma^*:S\Rightarrow_G^*w\}$.  For $A\in V$, put
$L_G(A):=\{w\in\Sigma^*:A\Rightarrow_G^*w\}$.
For a nonterminal $X\in V$, a terminal context $(u,v)$ is a
\emph{terminal reaching context of $X$ in $G$} if
$S\Rightarrow_G^*uXv$.
A nonterminal is \emph{reachable} if it occurs in a sentential form reachable
from $S$, and \emph{productive} if it derives a terminal word.  A nonterminal
is \emph{useless} if it is not both reachable and productive, and it is
\emph{nullable} if it derives $\lambda$.  A production $A\to B$ with
$A,B\in V$ is a \emph{unit production} (or \emph{unit rule}).  The
\emph{unit closure} is the reflexive-transitive closure of the relation on
nonterminals induced by unit productions.  A grammar is \emph{reduced} if
every nonterminal is reachable and productive.  For a grammar generating a
nonempty language, its \emph{trim} (or \emph{reduced part}) is the subgrammar
obtained by deleting all useless nonterminals together with their incident
productions.  A CFG is \emph{linear} if every right-hand side contains at
most one nonterminal.  By \emph{binarization} we mean replacement of longer
right-hand sides by equivalent chains of binary productions using fresh
nonterminals.
We write $\mathrm{CFL}$ for the class of context-free languages.  For any
finite representation $r$, $|r|$ denotes the length of a fixed reasonable
encoding; for a grammar this is polynomially equivalent to the usual symbol
count of its productions.  Following Yoshinaka's fixed-window analysis~\cite{yoshinaka2008}, for a reduced grammar $G=(V,\Sigma,P,S)$ we use the \emph{thickness} parameter
\[
 \tau_G:=\max_{A\in V}\min\{\,|w|:A\Rightarrow_G^*w,\ w\in\Sigma^*\,\}.
\]

\subsection{Learning from positive data}\label{sec:learning}

We use Gold's explanatory model~\cite{gold1967}.  Let $H_0$ denote a fixed
CFG with one start symbol and no productions, so $L(H_0)=\emptyset$.  We use
texts only for nonempty targets and represent the empty language separately by
$H_0$; equivalently one may adjoin the usual pause symbol.  For a nonempty language $L$, a \emph{text} is an
infinite sequence $T=(w_1,w_2,\ldots)$ over $L$ in which every word of $L$
occurs.  Write $T[n]=(w_1,\ldots,w_n)$ and
$K_n:=\{w_1,\ldots,w_n\}$.  A learner $\mathcal A$
\emph{identifies a nonempty $L$ in the limit} if for every text $T$ for $L$
there are $n_0$ and a grammar $H$ such that $\mathcal A(T[n])=H$ for every
$n\ge n_0$ and $L(H)=L$.
Thus the hypothesis representation itself eventually stabilizes.  A
sequential learner is \emph{conservative} if it changes its current hypothesis
only when the newly presented positive datum is not generated by that
hypothesis.

We distinguish this sequential learner from a set-driven finite-sample
operator $\mathcal B$, where $\mathcal B(K)$ depends only on the finite set
$K$, not on presentation order or multiplicity.  We call $\mathcal B$
\emph{sample-consistent} if $K\subseteq L(\mathcal B(K))$ for every finite
sample $K$.  A finite set $C\subseteq L$ is a \emph{characteristic sample for
$\mathcal B$} if
$C\subseteq K\subseteq L$ implies $L(\mathcal B(K))=L$ for every finite
$K$.  This is the finite-sample component of de~la~Higuera's polynomial
grammatical-inference framework~\cite{higuera1997}.  We measure sample size by
$\|K\|:=\sum_{w\in K}(|w|+1)$%
\footnote{The $+1$ prevents a nonempty sample such as $K=\{\lambda\}$ from having size $0$; it also gives $|K|\le\|K\|$, so the number of examples is controlled by the same measure.},
so $\lambda$ contributes one unit; the quantitative results below combine
polynomial-time construction of $\mathcal B(K)$ with polynomial-size
characteristic data.

\subsection{Fixed finite-monoid typing}\label{sec:fixed-typing}

Fix a finite monoid $M$ and a homomorphism $h:\Sigma^*\to M$, represented
by the multiplication table of $M$ and the letter images $h(a)$ for
$a\in\Sigma$.  Thus $h(w)$ is computable by a left-to-right product in
$O(|w|)$ time; this is the sense in which the fixed homomorphism is
\emph{efficiently computable} below.

\begin{definition}[fixed-$h$ substitutability]\label{def:hsubst}
A language $L\subseteq\Sigma^*$ is \emph{$\sim_h$-substitutable} if for all
$x,y\in\Sigma^+$, equality $h(x)=h(y)$ together with
$\D_L(x)\cap\D_L(y)\ne\varnothing$ implies $\D_L(x)=\D_L(y)$.
\end{definition}
The empty word is not used as an internal substitutable fragment.  When
$\lambda\in L$, the reconstruction handles it separately through the start
symbol.

For a nonempty internal factor $x$, we call $h(x)$ its
\emph{yield type}.  We write $\Ccf{h}$ for the context-free
$\sim_h$-substitutable languages and define
$\Clin{h}:=\{L\in\Ccf{h}:L\text{ is generated by some linear CFG}\}$.
The polynomial data bound below is measured relative to a chosen linear CFG
representation of the target; the reconstruction hypothesis produced later
need not itself be linear.

\section{Fixed Recognizable Substitutability}
\label{sec:framework}

For the fixed homomorphism $h$, put
$\mathsf{RS}_h:=\{L\subseteq\Sigma^*:L\text{ is $\sim_h$-substitutable}\}$.
Thus $\Ccf{h}=\mathsf{RS}_h\cap\mathrm{CFL}$, while $\Clin{h}$ consists of
those languages in $\mathsf{RS}_h\cap\mathrm{CFL}$ that are generated by
some linear CFG.
For expressive comparisons only, let
$\mathsf{RS}:=\bigcup_h\mathsf{RS}_h$, where $h$ ranges over finite-monoid
homomorphisms with domain $\Sigma^*$.  All such class notation is relative to
the ambient alphabet.  When the alphabet must be displayed explicitly, we
write $\mathsf{RS}_h(\Gamma)$ and $\mathsf{RS}(\Gamma)$ for the analogous
classes over $\Gamma^*$; class inclusions are always understood
alphabetwise.  The learning results in this paper concern one fixed $h$.
For every nonempty alphabet, neither $\mathsf{RS}$ nor
$\mathsf{RS}\cap\mathrm{CFL}$ is identifiable from positive data.\footnote{By
Proposition~\ref{prop:regular-auto}, every regular language belongs to
$\mathsf{RS}$. Hence both $\mathsf{RS}$ and $\mathsf{RS}\cap\mathrm{CFL}$
contain every finite language over the ambient alphabet and also the infinite
regular language $\Sigma^*$. Thus both classes are superfinite, and Gold's
theorem~\cite{gold1967} applies.}  There is, however, a simple monotonicity
with respect to refinement of the typing bias.  If $h':\Sigma^*\to M'$ refines $h$ in the sense that
$\ker h'\subseteq\ker h$, then
$\mathsf{RS}_h\subseteq\mathsf{RS}_{h'}$: equality under $h'$ implies
equality under $h$, so every comparison required by $h'$ is already covered by
$h$-substitutability.  Thus finer typings permit a larger fixed-$h$ class by
requiring fewer pairs of factors to be compared.  This gain is not cost-free
in the quantitative bounds below, since refining the typing may increase the
number of typed nonterminals and hence the witness-set size.  In particular,
for two finite typings $h$ and $g$, the product homomorphism $h\times g$ has
$\ker(h\times g)=\ker h\cap\ker g$ and therefore refines both; hence
$\mathsf{RS}_h\cup\mathsf{RS}_g\subseteq\mathsf{RS}_{h\times g}$.  Thus
window information and an independent finite domain-specific feature can be
combined in a single fixed typing.

A second useful observation gives a simple sufficient condition on the typing.
If $\ker h$, restricted to $\Sigma^+$, refines $\equiv_L$, then%
\footnote{If one required $\ker h\subseteq\equiv_L$ on all of $\Sigma^*$, the condition would be too strong: for the trivial monoid, $\lambda$ and every nonempty word have the same type, so even the ordinarily substitutable language $a^+$ would be excluded because $\lambda\not\equiv_{a^+}a$.}
$L\in\mathsf{RS}_h$, since equal $h$-types are already syntactically congruent.
This is not necessary: Proposition~\ref{prop:linear-separator-example} below
gives a four-element fixed-$h$ example whose syntactic congruence has infinite
index whereas every finite-monoid kernel has finite index.  Thus fixed-$h$ substitutability is
strictly weaker than requiring $\ker h$ to refine the full syntactic congruence.

\begin{proposition}\label{prop:regular-auto}
A language is regular iff it has the form $h^{-1}(\mathsf{Acc})$ for some homomorphism
$h:\Sigma^*\to M$ into a finite monoid $M$ and some $\mathsf{Acc}\subseteq M$.  For every fixed
$h$ and $\mathsf{Acc}\subseteq M$, the language $h^{-1}(\mathsf{Acc})$ is
$\sim_h$-substitutable.  Thus every regular language belongs to
$\mathsf{RS}$.
\end{proposition}

\begin{proof}
The first statement is the standard finite-monoid characterization of regular
languages~\cite{eilenberg1974}.  If $L=h^{-1}(\mathsf{Acc})$ and $h(x)=h(y)$,
then $h(uxv)=h(uyv)$ for every $u,v$; hence $\D_L(x)=\D_L(y)$.
\end{proof}

\subsection{Classical special cases}

A language $L\subseteq\Sigma^*$ is \emph{Clark--Eyraud substitutable} if, for all $x,y\in\Sigma^+$,
$\D_L(x)\cap\D_L(y)\ne\varnothing$ implies $\D_L(x)=\D_L(y)$.
Under the nonempty-fragment convention of Definition~\ref{def:hsubst}, this
is exactly the trivial-monoid case.  In particular, it is Yoshinaka's
$(0,0)$ case; see also Clark~\cite{clark2013} for the same operational
convention.

\begin{proposition}[Clark--Eyraud special case]\label{prop:ce-special}
Every Clark--Eyraud substitutable language belongs to $\mathsf{RS}$.
\end{proposition}
\begin{proof}
Take the one-element monoid.  Then the type condition is vacuous and
Definition~\ref{def:hsubst} is ordinary substitutability.
\end{proof}

For $k,\ell\ge0$, write $\pref_k(w)$ and $\suf_\ell(w)$ for the prefix and
suffix of the indicated lengths whenever they exist, with
$\pref_0(w)=\suf_0(w)=\lambda$.

Yoshinaka's original definition~\cite{yoshinaka2008} writes the left
length-$k$ window as $v\in\Sigma^k$ and the right length-$\ell$ window as
$u\in\Sigma^\ell$.  To avoid collision with our generic context notation
$(u,v)$, we rename only these two window variables to $p$ and $q$ below.
Thus a language $L$ is \emph{$(k,\ell)$-substitutable} if for all
$p\in\Sigma^k$, $q\in\Sigma^\ell$ and strings
$x_1,x_2,y_1,y_2,z_1,z_2\in\Sigma^*$, with $py_1q,py_2q\ne\lambda$,
\[
\begin{split}
&x_1py_1qz_1\in L,\quad x_1py_2qz_1\in L,\quad
x_2py_1qz_2\in L\\
&\hspace{35mm}\Longrightarrow\quad x_2py_2qz_2\in L.
\end{split}
\]

\begin{proposition}\label{prop:yl-special}
For every fixed $k,\ell\ge0$ there is a finite monoid homomorphism
$h_{k,\ell}$ such that, for every language $L\subseteq\Sigma^*$,
$L$ is $(k,\ell)$-substitutable if and only if it is
$\sim_{h_{k,\ell}}$-substitutable.
\end{proposition}

\begin{proof}

Put $m_0:=\max\{1,k+\ell\}$.  Define
\[
h_{k,\ell}(w)=
\begin{cases}
 w, & |w|<m_0,\\
 \langle\pref_k(w),\suf_\ell(w)\rangle, & |w|\ge m_0,
\end{cases}
\]
where a one-bit tag is used so that the elements representing short words and
those representing long words form disjoint sets, and put
$M_{k,\ell}:=h_{k,\ell}(\Sigma^*)$.  For short words $x,y$ and tagged
elements in the image, define multiplication by
\[
\begin{aligned}
 x\cdot y&:=h_{k,\ell}(xy), &
 x\cdot\langle p,q\rangle&:=\langle\pref_k(xp),q\rangle,\\
 \langle p,q\rangle\cdot y&:=\langle p,\suf_\ell(qy)\rangle, &
 \langle p,q\rangle\cdot\langle p',q'\rangle&:=\langle p,q'\rangle.
\end{aligned}
\]
A direct check of the four cases gives
$h_{k,\ell}(xy)=h_{k,\ell}(x)\cdot h_{k,\ell}(y)$.  Since $h_{k,\ell}$ is
onto, associativity follows from associativity of concatenation, and
$h_{k,\ell}(\lambda)$ is the identity.  Thus $M_{k,\ell}$ is a finite
monoid and $h_{k,\ell}$ is a homomorphism.

Suppose $s,t\in\Sigma^+$, $h_{k,\ell}(s)=h_{k,\ell}(t)$, and
$\D_L(s)\cap\D_L(t)\ne\varnothing$.  If $|s|<m_0$, then $s=t$.
Otherwise $s,t$ are both long and have the same length-$k$ prefix $p$ and
length-$\ell$ suffix $q$, so write $s=py_1q$ and $t=py_2q$.
Let $(u_0,v_0)$ be a shared context of $s,t$, and let $(u',v')$ be any
context of $s$.  The three words
\[
u_0py_1qv_0,\qquad u_0py_2qv_0,\qquad u'py_1qv'
\]
belong to $L$.  Applying $(k,\ell)$-substitutability gives
$u'py_2qv'\in L$.  Hence $\D_L(s)\subseteq\D_L(t)$, and symmetry gives
equality.

Conversely, assume that $L$ is $\sim_{h_{k,\ell}}$-substitutable and
consider strings satisfying the premises of the $(k,\ell)$-substitutability
condition.  Put $s:=py_1q$ and $t:=py_2q$; by definition both are nonempty.
If $k+\ell>0$, then $|s|,|t|\ge k+\ell=m_0$ and both have type
$\langle p,q\rangle$.  If $k=\ell=0$, nonemptiness gives $|s|,|t|\ge1=m_0$,
and both have the unique nonempty-word type $\langle\lambda,\lambda\rangle$.
Thus $h_{k,\ell}(s)=h_{k,\ell}(t)$.  The first two premise words show that
$s$ and $t$ share context $(x_1,z_1)$, while the third gives
$(x_2,z_2)\in\D_L(s)$.  Fixed-$h_{k,\ell}$ substitutability therefore gives
$(x_2,z_2)\in\D_L(t)$, which is exactly the required fourth word.
\end{proof}

\phantomsection\label{crit:kl}
\paragraph{Fixed-window counterexample criterion.}
Let $k,\ell\ge0$ and put $m_0:=\max\{1,k+\ell\}$.  If
$s,t\in\Sigma^*$ satisfy $|s|,|t|\ge m_0$,
$\pref_k(s)=\pref_k(t)$, $\suf_\ell(s)=\suf_\ell(t)$,
$\D_L(s)\cap\D_L(t)\ne\varnothing$, and $\D_L(s)\ne\D_L(t)$,
then $L$ is not $(k,\ell)$-substitutable.  Indeed, the length and
prefix--suffix assumptions give $h_{k,\ell}(s)=h_{k,\ell}(t)$, whereas the
shared context together with unequal distributions shows that $L$ is not
$\sim_{h_{k,\ell}}$-substitutable.  Proposition~\ref{prop:yl-special}
therefore gives the claim.

Consequently the fixed $(k,\ell)$ class is exactly
$\mathsf{RS}_{h_{k,\ell}}$.  Section~\ref{sec:fixed-window-thick} verifies
that the present reconstruction also recovers Yoshinaka's
grammar-size-and-thickness finite-sample scale on this special case.  For
$k=\ell=0$, the homomorphism $h_{0,0}$ separates $\lambda$ from the nonempty
words, so under Definition~\ref{def:hsubst} it induces the same
substitutability condition as the one-element trivial monoid.  We write
$\mathsf{KL}$ for the union over all fixed $k,\ell$ on the current alphabet,
and $\mathsf{KL}(\Gamma)$ when the alphabet $\Gamma$ is displayed explicitly.

\subsection{Main theorem}

\begin{theorem}[Fixed-$h$ learning theorem]\label{thm:main}
Fix an efficiently computable homomorphism $h:\Sigma^*\to M$ into a finite
monoid.  There are a set-driven reconstruction operator $\mathcal B_h$ and
a conservative sequential learner $\mathcal A_h$ such that:
\begin{enumerate}[label=\textup{(\roman*)},nosep]
\item $\mathcal B_h$ is sample-consistent, and $\mathcal B_h(K)$ is
constructible in time polynomial in $\|K\|$ for every finite positive
sample $K$;
\item every $L\in\Ccf{h}$ has a finite characteristic
sample $C_L\subseteq L$ for $\mathcal B_h$; and
\item $\mathcal A_h$ identifies every nonempty $L\in\Ccf{h}$ in Gold's
sense, and each update is polynomial-time in $\|K_n\|$, where $K_n$ is the
positive-example set seen up to that stage.
\end{enumerate}
For every fixed $(k,\ell)$, characteristic data for targets in
$\Ccf{h_{k,\ell}}$---equivalently, for context-free $(k,\ell)$-substitutable
targets---can moreover be chosen with encoded size polynomial in reduced
target CFG size and thickness.  For the linear subclass $\Clin{h}$ under an
arbitrary fixed $h$, the characteristic samples can be chosen with encoded
size polynomial in the size of a linear grammar for the target.
\end{theorem}

Sections~\ref{sec:algo}--\ref{sec:complexity} prove the qualitative theorem;
Sections~\ref{sec:fixed-window-thick} and~\ref{sec:linear} establish the two
quantitative bounds.  For the linear statement, Section~\ref{sec:linear}
derives direct polynomial bounds on the canonical yields and canonical
contexts used by the witness construction.

There is no conflict with Gold's standard superfinite obstruction: a class
containing all finite languages together with an infinite language is not
identifiable from positive data.  Let the alphabet be nonempty and fix any
finite $h$.  Choose distinct nonempty strings $x,y$ with $h(x)=h(y)$, and
then choose $r$ with $|r|>\max\{|x|,|y|\}$.  Then
$rx\notin\{x,y\}$ and $ry\notin\{x,y,rx\}$: the first exclusion follows
from length, while $rx\ne ry$ follows from left cancellation and $x\ne y$.
The finite language $\{x,y,rx\}$ is not $\sim_h$-substitutable.  Indeed,
$x$ and $y$ share the empty context, while $(r,\lambda)$ is a context of
$x$ but not of $y$.  Hence neither $\mathsf{RS}_h$ nor
$\Ccf{h}=\mathsf{RS}_h\cap\mathrm{CFL}$ contains all finite languages.

\section{Reconstruction Operator and Sequential Learner}
\label{sec:algo}


\subsection{Finite-sample reconstruction}
For a finite sample $K\subseteq\Sigma^*$, let
$\widehat V(K):=\{[x:u,v]:uxv\in K,\ x\in\Sigma^+\}$.
These are syntactic triples: $[x:u,v]=[x':u',v']$ iff
$x=x'$, $u=u'$, and $v=v'$; different observed factors or contexts therefore
give different hypothesis nonterminals.
The grammar $\mathcal B_h(K)=\widehat G(K)$ has nonterminal set
$\widehat V(K)\cup\{\widehat S\}$ and start symbol $\widehat S$.  It has the
following productions whenever every displayed bracketed nonterminal belongs
to $\widehat V(K)$:
\[
\begin{array}{rll}
(R1)&[xy:u,v]\to[x:u,yv][y:ux,v],&x,y\in\Sigma^+,\\[1mm]
(R2)&[x:u,v]\to[x:u',v'],&uxv,u'xv'\in K,\\[1mm]
(R3)&[x:u,v]\to[x':u,v],&h(x)=h(x'),\\[1mm]
(R4)&[a:u,v]\to a,&a\in\Sigma,\\[1mm]
(R5)&\widehat S\to[w:\lambda,\lambda],&w\in K\cap\Sigma^+.
\end{array}
\]
If $\lambda\in K$, Rule~(R5) also contains $\widehat S\to\lambda$.
Intuitively, $[x:u,v]$ represents possible replacements for the observed
factor $x$ in the context $(u,v)$.  For a $\sim_h$-substitutable target, the
soundness invariant below ensures that every derived word has the same
$h$-type as $x$ and is distributionally interchangeable with it.
Rule~(R2) transports the same observed factor between contexts in which
it has been observed; no condition on the $h$-types of those surrounding
contexts is needed.  Rule~(R3) changes the internal representative within a
fixed yield type, preserving the yield-type invariant used by soundness.

\begin{lemma}[sample consistency]\label{lem:sample-consistency}
For every finite sample $K$, $K\subseteq L(\mathcal B_h(K))$.
\end{lemma}
\begin{proof}
If $w=a_1\cdots a_m\in K\cap\Sigma^+$, Rule~(R5) enters
$[w:\lambda,\lambda]$.  Repeated use of Rule~(R1) along any binary
factorization of this occurrence of $w$ reaches the one-letter factors, and
every intermediate observed nonterminal exists because it is a factor of the
same sample word with its induced context.  Rule~(R4) then emits
$a_1\cdots a_m$.  If $\lambda\in K$, the start rule emits it directly.
\end{proof}

\begin{theorem}[soundness]\label{thm:soundness}
If $K\subseteq L$ and $L$ is $\sim_h$-substitutable, then
$L(\widehat G(K))\subseteq L$.
\end{theorem}

\begin{proof}
Induct on the height of a derivation tree rooted at $[x:u,v]$ and prove simultaneously that
\begin{equation}\label{eq:soundness-invariant}
 [x:u,v]\Rightarrow_{\widehat G(K)}^*w
 \quad\Longrightarrow\quad
 uwv\in L\ \text{ and }\ h(w)=h(x).
\end{equation}
Rules~(R1)--(R4) contain no erasing production, and every lexical rule emits
one terminal.  Hence every terminal word derived from a non-start hypothesis
nonterminal is nonempty; in particular, the words $w,w_1,w_2$ to which
Definition~\ref{def:hsubst} is applied below lie in $\Sigma^+$.
Rule~(R4) is immediate.  For Rule~(R3), the induction hypothesis gives
$h(w)=h(x')=h(x)$ and already gives $uwv\in L$.  For Rule~(R2), the
induction hypothesis gives $u'wv'\in L$ and $h(w)=h(x)$.  Since
$[x:u',v']$ and $[x:u,v]$ belong to $\widehat V(K)$, its definition gives
$u'xv',uxv\in K\subseteq L$.  Thus $x$ and $w$ share the context
$(u',v')$, so substitutability gives $\D_L(x)=\D_L(w)$ and hence
$uwv\in L$.

For Rule~(R1), write the derived word as $w_1w_2$.  The two induction
hypotheses give
\[
 uw_1yv\in L,\qquad uxw_2v\in L,\qquad
 h(w_1)=h(x),\qquad h(w_2)=h(y).
\]
Since $uxyv\in K$, the strings $x$ and $w_1$ share context $(u,yv)$ and
have the same type; hence their distributions agree.  By the second induction
hypothesis, $(u,w_2v)\in\D_L(x)$.  Applying the distribution equality to this
context gives $uw_1w_2v\in L$; also $h(w_1w_2)=h(xy)$ by the induction hypotheses.
This proves~\eqref{eq:soundness-invariant}.  Start rules then give the claim for $L(\widehat G(K))$;
$\lambda$ is present only when
$\lambda\in K\subseteq L$.
\end{proof}

\subsection{The sequential learner}

In general, the batch operator itself need not stabilize syntactically when
recomputed after every positive example.\footnote{For instance, for
$L=\{a^ncb^n:n\ge 0\}$, if increasingly long positive examples are added to
the accumulated sample, each new word contributes new factor/context triples
such as $[a^ncb^n:\lambda,\lambda]$.}  We therefore use the standard
conservative wrapper, starting from the fixed empty grammar $H_0$ defined in
Section~\ref{sec:learning}.  After reading $w_n$, put
\[
 H_n:=
 \begin{cases}
 H_{n-1},&w_n\in L(H_{n-1}),\\
 \mathcal B_h(K_n),&w_n\notin L(H_{n-1}),
 \end{cases}
 \qquad
 \mathcal A_h(T[n]):=H_n.
\]
The finite-sample characteristic-set viewpoint follows the classical
substitutability literature; cf. Clark--Eyraud~\cite{clark-eyraud2007}.
Here we place the batch constructor inside a conservative sequential wrapper:
every rebuild uses the entire accumulated set $K_n$, and a rebuild occurs only
when the current hypothesis rejects the next positive example.  Consequently,
once a characteristic sample for $\mathcal B_h$ is contained in $K_n$, any later
rebuild falls directly under Theorem~\ref{thm:reconstruction-fixed-h}.

\section{Completeness via Yield-Typed Refinement and Finite Witnesses}
\label{sec:typing}


\subsection{Target refinement and finite witnesses}

For the completeness proof only, we refine a target grammar by yield type and
show that the hypothesis can mimic finitely witnessed typed derivations.  We use
the following normal form as a proof device.

\begin{definition}[start-separated binary normal form]
A CFG $G=(V,\Sigma,P,S_0)$ is in \emph{start-separated binary normal form}
(SSBNF) if $S_0$ never occurs on a right-hand side%
\footnote{This isolates whole-sentence entry from internal nonterminals, simplifying the typed refinement and the special treatment of $\lambda$.},
every start rule is $S_0\to A$ or $S_0\to\lambda$, and every non-start rule
is either $A\to a$ or $A\to BC$.
\end{definition}

Every nonempty context-free language has an equivalent reduced SSBNF grammar by
the standard start-separation, useless-symbol elimination, $\lambda$/unit elimination
away from the start, and binarization procedures; the case of the empty
target language is immediate.

Let $G=(V,\Sigma,P,S_0)$ be a reduced SSBNF grammar for
$L\in\Ccf{h}$.  Its \emph{full yield-typed refinement} has start symbol
$\widetilde S$ and non-start symbols $A_\mu$ with
$A\in V\setminus\{S_0\}$ and $\mu\in M$.  Here and below, juxtaposition of type indices denotes multiplication in
$M$; thus $\mu\nu=\mu\cdot\nu$.  The rules are
\[
\begin{array}{ll}
\widetilde S\to A_\mu & (S_0\to A\in P,\ \mu\in M),\\
\widetilde S\to\lambda & (S_0\to\lambda\in P),\\
A_{h(a)}\to a & (A\to a\in P),\\
A_{\mu\nu}\to B_\mu C_\nu & (A\to BC\in P,\ \mu,\nu\in M).
\end{array}
\]
The \emph{typed refinement}
$\widetilde G=(\widetilde V,\Sigma,\widetilde P,\widetilde S)$ is the
reduced part of this grammar.  By convention, every symbol of
$\widetilde G$ is therefore reachable and productive.  The yield-type index
is needed for Rule~(R3) in the completeness proof: without it, for a target
rule $A\to BC$, the canonical yield $\omega(A)$ and the concatenation
$\omega(B)\omega(C)$ need not have the same $h$-value, so the required
Rule~(R3) transition need not exist.  Splitting by yield type makes this
equality an invariant.  By contrast, no types are attached to the
surrounding left or right contexts, because the reconstruction operator never
tests such types.

\begin{proposition}[typed lifting and yield invariant]\label{prop:typed-core}
The typed refinement satisfies $L(\widetilde G)=L(G)$.  Moreover, if $X=A_\mu$ is a non-start symbol of $\widetilde G$, then
$X\Rightarrow_{\widetilde G}^*w$ implies $h(w)=\mu$.
\end{proposition}
\begin{proof}
Erase type annotations to map every typed derivation to a derivation of $G$.
Conversely, annotate each non-start node of a successful derivation tree of
$G$ by the $h$-type of its terminal yield.  The terminal and binary rules
above are exactly the bottom-up propagation laws for those types, so the
annotated tree is a derivation of the full refinement; all symbols on a
successful tree survive trimming.  The invariant follows by induction on a
typed derivation.
\end{proof}

Fix an order on $\Sigma$ and use the induced \emph{shortlex} order:
shorter strings come first, with lexicographic order breaking equal-length
ties.  For every non-start symbol $X$ of $\widetilde G$, let $\omega(X)$
be its shortlex-minimal terminal yield, called the \emph{canonical yield} of $X$.  Choose $\chi(X)=(u_X,v_X)$ among the terminal reaching contexts of $X$ in
$\widetilde G$ so as to minimize $|u_X|+|v_X|$, breaking ties by the
shortlex orders of $u_X$ and then $v_X$; call $\chi(X)$ the \emph{canonical context} of $X$.  Such a terminal reaching context
exists because $\widetilde G$ is reduced:
reachability places $X$ in a reachable sentential form, and productivity of
the surrounding nonterminals lets all symbols to its left and right be
expanded to terminal strings.  Proposition~\ref{prop:typed-core} gives
$h(\omega(A_\mu))=\mu$ for every typed non-start symbol.

Define the finite witness set
\begin{equation}\label{eq:W}
\begin{split}
\WitnessSet(\widetilde G):={}&
 \{u_X\omega(X)v_X:X\in\widetilde V\setminus\{\widetilde S\}\}\\
&\cup\{u_Xav_X:X\to a\in\widetilde P\}\\
&\cup\{u_X\omega(Y)\omega(Z)v_X:X\to YZ\in\widetilde P\}\\
&\cup\bigl(\{\lambda\}\text{ if }\lambda\in L\bigr).
\end{split}
\end{equation}
Every displayed word belongs to $L$, and the set is finite.  In particular,
for a binary rule $X\to YZ$ the corresponding witness is justified by
$\widetilde S\Rightarrow^*u_XXv_X\Rightarrow u_XYZv_X
\Rightarrow^*u_X\omega(Y)\omega(Z)v_X$; anchors and terminal-rule
witnesses are analogous.  We call the
elements of $\WitnessSet(\widetilde G)$ the \emph{canonical positive witnesses}.
Here the word $u_X\omega(X)v_X$ is the \emph{anchor} of $X$, witnessing one canonical
occurrence of $X$, while the other words serve the distinct role of terminal or binary
\emph{rule witnesses}.


\subsection{Exact reconstruction and Gold convergence}

\begin{theorem}[completeness from the finite witnesses]\label{thm:complete}
Let $G$ be any reduced SSBNF grammar for a nonempty language $L$, and let
$\widetilde G$ be its typed refinement with respect to the fixed
homomorphism $h$.  If a finite set $K$ satisfies
$\WitnessSet(\widetilde G)\subseteq K$, then $L\subseteq L(\widehat G(K))$.
\end{theorem}

The statement is purely combinatorial: it uses the typed grammar and
$\WitnessSet(\widetilde G)\subseteq K$, but neither substitutability nor the
additional assumption $K\subseteq L$.

\begin{proof}
For each non-start $X$ write $\widehat X:=[\omega(X):u_X,v_X]$.
The anchor word in~\eqref{eq:W} ensures that $\widehat X$ exists.  We prove
by induction on a typed derivation $X\Rightarrow^*w$ that
$\widehat X\Rightarrow_{\widehat G(K)}^*w$.

If the first rule is $X\to a$, the witness $u_Xav_X\in K$ gives
$[a:u_X,v_X]$.  If $\omega(X)\ne a$, Rule~(R3) changes $\widehat X$ to
$[a:u_X,v_X]$ because $h(\omega(X))=h(a)$; if $\omega(X)=a$, this step is
unnecessary.  Rule~(R4) then yields $a$.

Suppose the first rule is $X=A_{\mu\nu}\to B_\mu C_\nu$; put
$Y:=B_\mu$ and $Z:=C_\nu$.
The rule witness in~\eqref{eq:W} gives the three observed nonterminals
needed to apply Rule~(R1), yielding
\begin{equation}\label{eq:binary-reconstruction}
 [\omega(Y)\omega(Z):u_X,v_X]
 \to[\omega(Y):u_X,\omega(Z)v_X]
      [\omega(Z):u_X\omega(Y),v_X].
\end{equation}
Because $h(\omega(X))=\mu\nu=h(\omega(Y)\omega(Z))$, Rule~(R3), if
necessary, changes $\widehat X$ to the left-hand side
of~\eqref{eq:binary-reconstruction}; when $\omega(X)=\omega(Y)\omega(Z)$,
this step is unnecessary.  Rule~(R2) then transports
the two children to $\widehat Y$ and $\widehat Z$, which are determined by
the canonical contexts of $Y$ and $Z$.  The induction hypotheses finish the two subderivations.

Finally, every nonempty $w\in L$ has a lifted start derivation
$\widetilde S\to X\Rightarrow^*w$.  The start rule itself shows that
$(\lambda,\lambda)$ is a terminal reaching context of $X$; its total length is $0$,
which is minimal, so the canonical choice is
$\chi(X)=(\lambda,\lambda)$.  Its anchor $\omega(X)$ belongs to $K$, so
Rule~(R5) starts the preceding derivation.  The case $w=\lambda$ is immediate
from~\eqref{eq:W} and Rule~(R5).
\end{proof}

\begin{theorem}[exact finite-sample reconstruction]
\label{thm:reconstruction-fixed-h}
Let $L\in\Ccf{h}$ be nonempty, let $G$ be a reduced SSBNF grammar for $L$,
and let $\widetilde G$ be its typed refinement.  If
$\WitnessSet(\widetilde G)\subseteq K\subseteq L$, then $L(\mathcal B_h(K))=L$.
\end{theorem}
\begin{proof}
Combine Theorems~\ref{thm:soundness} and~\ref{thm:complete}.
\end{proof}

Thus, for every nonempty target representation $G$, the set
$\WitnessSet(\widetilde G)$ is a characteristic sample for $L(G)$
with respect to $\mathcal B_h$ in the sense of Section~\ref{sec:learning}.

Theorem~\ref{thm:reconstruction-fixed-h} covers every nonempty target,
including $\{\lambda\}$.  For the remaining endpoint $L=\emptyset$, take
$C_L=\emptyset$; then $\mathcal B_h(\emptyset)$ has no start production.
Together these cases establish item~(ii) of Theorem~\ref{thm:main}.

\begin{corollary}[Gold identification]\label{cor:ilt}
For a fixed homomorphism $h:\Sigma^*\to M$ into a finite monoid $M$, the learner
$\mathcal A_h$ identifies every nonempty language in $\Ccf{h}$ from positive data.  More precisely, for any reduced SSBNF target
grammar $G$ and its typed refinement $\widetilde G$, once
$\WitnessSet(\widetilde G)\subseteq K_{n_0}$, among stages $n\ge n_0$ at most one hypothesis change
can occur.
\end{corollary}

\begin{proof}
Under our text convention the Gold statement concerns nonempty targets;
$H_0$ is the initial empty hypothesis.  Let $L$ be nonempty and let $T$ be a
text for $L$.  Since $w_1\notin L(H_0)=\emptyset$, the learner rebuilds at
stage $1$.  Choose $n_0$ with
$\WitnessSet(\widetilde G)\subseteq K_{n_0}$.  If the learner rebuilds at
some stage $n\ge n_0$, Theorem~\ref{thm:reconstruction-fixed-h} makes the
new hypothesis exactly $L$, after which no positive example can trigger
another change.  In this case that rebuild is the last hypothesis change.
Suppose instead that no rebuild occurs at any stage $n\ge n_0$.  Soundness
gives $L(H_{n_0})\subseteq L$.  Moreover, $K_n\subseteq L(H_n)$ for every
$n$: a rebuild is sample-consistent, while at a non-rebuild step the newly
presented datum is already generated by the unchanged hypothesis.  Hence a
strict inclusion $L(H_{n_0})\subsetneq L$ would leave some
$w^\star\in L\setminus L(H_{n_0})$ that has not yet occurred by stage
$n_0$.  Since every target word eventually occurs in the text, $w^\star$
must occur later and force a rebuild, a contradiction.  Thus
$L(H_{n_0})=L$; in this case no hypothesis change after stage $n_0$ is
needed.
\end{proof}

\section{Complexity of the Reconstruction Step}
\label{sec:complexity}

\begin{theorem}\label{thm:poly-build}
For fixed $h$, $\mathcal B_h(K)$ is constructible in time polynomial in
$\|K\|$.
\end{theorem}
\begin{proof}
Put $n_K:=\|K\|$.  Across all sample words there are $O(n_K^2)$
factorizations $uxv$ with $x\ne\lambda$, so
$|\widehat V(K)|=O(n_K^2)$.  All substring $h$-values can also be
precomputed and cached using $O(n_K^2)$ monoid multiplications.

Rule~(R1) has $O(n_K^3)$ candidates, whereas Rules~(R2) and~(R3) use pairs
of observed nonterminals and therefore have at most $O(n_K^4)$ candidates.
For fixed $h$, each type test is constant-time.  Hence the total number of
productions generated is at most $O(n_K^4)$.

Each nonterminal occurring in a production explicitly stores a factor and its
context from a sample word, so the encoding length of each production is
$O(n_K)$.  Therefore the total encoding length of the output grammar is
$O(n_K^4)\cdot O(n_K)=O(n_K^5)$.  If factors and contexts are represented
by precomputed occurrence identifiers, the equality tests used while scanning
candidate tuples are constant-time; even with direct comparison of their
$O(n_K)$-length encodings, scanning and checking all $O(n_K^4)$ candidates
takes at most $O(n_K^5)$ time.  Explicitly writing the output grammar also
takes at most $O(n_K^5)$ time.  Thus $\mathcal B_h(K)$ can be constructed
explicitly in $O(n_K^5)$ time.
\end{proof}

\begin{corollary}\label{cor:poly-update}
At each stage the conservative learner $\mathcal A_h$ can process the next
positive example in time polynomial in the encoded size of the data seen so
far.
\end{corollary}

\begin{proof}
CFG membership is polynomial in the grammar and input sizes, and every current
hypothesis has polynomial size by Theorem~\ref{thm:poly-build}; a triggered
rebuild is polynomial by the same theorem.
\end{proof}

Lemma~\ref{lem:sample-consistency}, Theorems~\ref{thm:reconstruction-fixed-h}
and~\ref{thm:poly-build}, together with Corollaries~\ref{cor:ilt}
and~\ref{cor:poly-update}, prove the qualitative part of Theorem~\ref{thm:main}.
A grammar-size-only data bound is impossible for arbitrary CFG representations:
the SSBNF family $S_0\to A_n$, $A_0\to a$,
$A_i\to A_{i-1}A_{i-1}$ has $O(n)$ production-symbol count, and hence encoded
size polynomial in $n$ under our convention, but generates only $a^{2^n}$.
More generally, every singleton language is $\sim_h$-substitutable for every
$h$: if two nonempty factors share a context inside its unique word, then
cancellation in the free monoid makes the factors
equal.  Thus this obstruction already occurs inside every fixed-$h$ target
class.  Every characteristic sample must contain a word of length $2^n$, which is
superpolynomial in the encoded grammar size.  The next section therefore uses
the thickness parameter used in Yoshinaka's fixed-window analysis; Section~\ref{sec:linear} later obtains a grammar-size-only
bound for linear targets under arbitrary fixed $h$.


\section{Fixed-Window Thick Data}
\label{sec:fixed-window-thick}

By Proposition~\ref{prop:yl-special}, the fixed-window case is exactly a
particular fixed finite-monoid typing.  The purpose of this section is therefore
a compatibility result.  Yoshinaka's Theorem~1, with Lemmas~6 and~7 separating the characteristic-set cardinality and description-size bounds, establishes a grammar-size-and-thickness characteristic-data scale~\cite{yoshinaka2008}.  We verify that the present reconstruction retains that same scale, rather than claiming a new learnability result for the fixed-window class.  For $h_{k,\ell}$, a
typed nonterminal $A_\mu$ need preserve only a bounded prefix and suffix,
which controls the additional witness cost.

Fix $k,\ell\ge0$, put $r:=k+\ell$, and let $h_{k,\ell}$ be the morphism of
Proposition~\ref{prop:yl-special}.  Let $G=(V,\Sigma,P,S_0)$ be a reduced
SSBNF grammar and let $\widetilde G$ be its reduced yield-typed refinement for
$h_{k,\ell}$.  Write
\[
 N:=|V\setminus\{S_0\}|,\qquad
 N_t:=|\widetilde V\setminus\{\widetilde S\}|,\qquad
 B_{k,\ell}(G):=
 \begin{cases}
   \tau_G, & r=0,\\[1mm]
   r+(2r-1)N\tau_G, & r>0.
 \end{cases}
\]

\begin{lemma}[bounded fixed-window typed yields]
\label{lem:window-typed-yield}
For every non-start typed symbol $A_\mu$ of $\widetilde G$,
$|\omega(A_\mu)|\le B_{k,\ell}(G)$.
\end{lemma}

\begin{proof}
See Appendix~\ref{app:fixed-window-bounds}.
\end{proof}

\begin{lemma}[bounded fixed-window contexts and witnesses]
\label{lem:window-context}
Every non-start typed symbol $X$ of $\widetilde G$ has a terminal
reaching context satisfying
\[
 |u_X|+|v_X|\le N_t B_{k,\ell}(G),
\]
and every word in $\WitnessSet(\widetilde G)$ has length at most
$(N_t+2)B_{k,\ell}(G)+1$.
\end{lemma}

\begin{proof}
See Appendix~\ref{app:fixed-window-bounds}.
\end{proof}

\begin{theorem}[polynomial thick data for fixed windows]
\label{thm:window-thick}
For every fixed pair $(k,\ell)$ and ambient alphabet $\Sigma$, there is a
polynomial $Q_{k,\ell}$ such that the following holds.  Let $G$ be a reduced
SSBNF grammar with $L(G)\in\Ccf{h_{k,\ell}}$, and let $\widetilde G$ be its
yield-typed refinement.  Then
\[
 \|\WitnessSet(\widetilde G)\|
 \le Q_{k,\ell}(|G|,\tau_G).
\]
Consequently $\WitnessSet(\widetilde G)$ is a characteristic sample for
$\mathcal B_{h_{k,\ell}}$ whose $\|\cdot\|$-size is polynomial in
$|G|$ and $\tau_G$.
\end{theorem}

\begin{proof}
The fixed monoid $M_{k,\ell}$ has constant size for fixed $k,\ell$ and
$\Sigma$.  Hence
$N_t\le N|M_{k,\ell}|$.  Each terminal production of $G$ yields at most one
productive typed terminal production, and each binary production yields at
most $|M_{k,\ell}|^2$ typed copies.  Therefore
$|\widetilde V|+|\widetilde P|=O_{k,\ell}(|G|)$.
The witness set contains at most one anchor per typed nonterminal, one witness
per typed terminal or binary production, and possibly $\lambda$.  Lemma
\ref{lem:window-context} bounds every witness length by a polynomial in
$|G|$ and $\tau_G$.  Multiplying the polynomial number of witnesses by this
length bound proves the claim.  Characteristicity is Theorem~\ref{thm:reconstruction-fixed-h}.
\end{proof}

The theorem above is stated for reduced SSBNF representations.  The next
normalization lemma transfers it to arbitrary reduced CFG representations.

\begin{proposition}[polynomial thickness-preserving SSBNF normalization]
\label{prop:thick-ssbnf-normal}
There are fixed polynomials $P_1$ and $P_2$ such that every reduced CFG
$R$ generating a nonempty language can be converted in polynomial time into
an equivalent reduced SSBNF grammar $G$ satisfying
\[
 |G|\le P_1(|R|),\qquad
 \tau_G\le P_2(|R|,\tau_R+1).
\]
\end{proposition}

The proof is given in Appendix~\ref{app:thick-ssbnf}; terminal isolation and
binarization are performed before non-start $\lambda$-elimination so that both
grammar size and shortest nonempty yields remain polynomially bounded.

\begin{corollary}[fixed-window complexity transfer]
\label{cor:window-transfer}
For every fixed $k,\ell\ge0$ and finite alphabet $\Sigma$, every language
$L\in\Ccf{h_{k,\ell}}$ represented by a reduced CFG $R$---equivalently,
every context-free $(k,\ell)$-substitutable target represented by such an
$R$---has characteristic data for $\mathcal B_{h_{k,\ell}}$ whose encoded
size is polynomial in $|R|$ and $\tau_R$.
\end{corollary}

\begin{proof}
Apply Proposition~\ref{prop:thick-ssbnf-normal} and then
Theorem~\ref{thm:window-thick}; the equivalence is
Proposition~\ref{prop:yl-special}.
\end{proof}

Thus, on the exactly corresponding fixed-window classes, the present batch
reconstruction recovers the same grammar-size-and-thickness finite-sample scale established by Yoshinaka's Theorem~1 and Lemmas~6--7~\cite{yoshinaka2008}.  This is a compatibility check for the
general reconstruction framework, not a new fixed-window learnability claim.
The polynomial depends on the fixed window and alphabet; no uniform bound in
variable $k,\ell$ is asserted.  Stabilization of the conservative sequential
learner is the separate Gold result of Corollary~\ref{cor:ilt}.

\section{The Linear Subclass}
\label{sec:linear}

Linear languages are exactly the languages accepted by one-turn pushdown
automata~\cite{ginsburg-spanier1966}.  The normalization below produces the
grammatical counterpart of the single-turn structure: every successful
derivation has one continuing nonterminal path.

\begin{proposition}[linear-spine SSBNF normalization]\label{prop:linear-normal}
There is a fixed polynomial $P_{\mathrm{lin}}$ such that every linear CFG $G$
generating a nonempty language can be converted in polynomial time into an equivalent reduced
SSBNF grammar $G_0=(V_0,\Sigma,P_0,S_0)$ with
$|G_0|\le P_{\mathrm{lin}}(|G|)$ in which every binary
production has exactly one fresh terminal-wrapper child $W_a$ with
$W_a\to a$ and exactly one other nonterminal child.
\end{proposition}

A complete construction is given in
Appendix~\ref{app:linear-normalization}.  We call the resulting grammar
\emph{linear-spine SSBNF}.  Let $\widetilde G_0$ be its typed refinement and put
$n_t:=|\widetilde V_0\setminus\{\widetilde S\}|$.  Since each underlying
non-start symbol has at most $|M|$ typed copies,
$n_t\le |V_0\setminus\{S_0\}|\,|M|$; for fixed $h$ this is polynomial in
$|G|$.  The spine and reaching-spine terminology used in the short-witness
proof is deferred to Appendix~\ref{app:linear-short-proof}.

\begin{lemma}[short canonical witnesses]\label{lem:linear-short}
For every typed non-start symbol $X$ of $\widetilde G_0$,
\[
 |\omega(X)|\le n_t,
 \qquad
 |u_X|+|v_X|\le 2n_t.
\]
Consequently every word in $\WitnessSet(\widetilde G_0)$ has length $O(n_t)$.
\end{lemma}

\begin{theorem}[polynomial characteristic data for linear targets]
\label{thm:linear-poly}
For fixed $h$, every language in $\Clin{h}$ represented by a linear CFG has a
characteristic sample for $\mathcal B_h$ whose $\|\cdot\|$-size is polynomial
in the target representation.
\end{theorem}

\begin{proof}
For $L=\emptyset$ and $L=\{\lambda\}$ use the characteristic samples from
Section~\ref{sec:typing}.  Otherwise apply
Proposition~\ref{prop:linear-normal} and Lemma~\ref{lem:linear-short}.
There are $O(|\widetilde V_0|+|\widetilde P_0|+1)$ canonical witnesses, each
of length $O(n_t)$, while $n_t$, $|\widetilde V_0|$, and
$|\widetilde P_0|$ are polynomial in the linear target representation for
fixed $h$.  Characteristicity follows from
Theorem~\ref{thm:reconstruction-fixed-h}.
\end{proof}

Together with Theorem~\ref{thm:poly-build}, this gives polynomial-time batch
reconstruction from arbitrary finite samples and polynomial characteristic data
relative to a chosen linear CFG representation.  We state this
representation-sensitive conclusion directly rather than label the sequential
learner itself as a de~la~Higuera polynomial time-and-data learner, since